\section{Related Work}
\label{sec:related}

The \Apeiron{} Framework occupies a unique position at the intersection of several active research traditions. In this section, we situate our contribution relative to foundational AI architectures, formal verification, non-anthropocentric epistemology, and the broader tradition of research programme papers in artificial intelligence.

\subsection{Foundational AI Architectures}
The ambition to build AI systems on rigorous mathematical foundations has a long history. Schmidhuber's \textit{G\"odel Machine} \cite{schmidhuber2007} proposed a self-referential agent capable of rewriting its own code, but it remains a theoretical construct without a layered ontology. Hutter's \textit{AIXI} \cite{hutter2005} provides a formal theory of optimal reinforcement learning based on Solomonoff induction, yet it assumes a fixed, human-defined reward signal and does not address the genesis of categories. More recently, \textit{World Models} \cite{ha2018} and \textit{Dreamer} \cite{hafner2019} have demonstrated that agents can learn compact latent representations of their environments; however, these models are still trained on human-curated reward functions and operate within a predefined ontology.

\subsection{Category Theory in AI}
The application of category theory to machine learning has gained traction, particularly through the lens of \textit{applied category theory} \cite{fong2019}. Frameworks such as \textit{DisCoPy} \cite{delpeuch2021} use monoidal categories to compose neural network layers, while \textit{categorical deep learning} \cite{gavranovic2024} explores functorial semantics of architectures. However, these approaches remain focused on representing existing human-designed architectures rather than enabling autonomous ontological creation. The \Apeiron{} Framework extends categorical methods beyond representation to the active construction of new categories (Layer~9) and the reconciliation of incommensurable ontologies (Layer~12).

\subsection{Formal Verification and Explainable AI}
The need for verifiable AI has spurred developments in \textit{neuro-symbolic verification} and \textit{certified learning}. Systems like \textit{AlphaVerus} \cite{wang2023} and \textit{VeriDeep} \cite{chen2024} provide guarantees for specific neural network properties, but they operate post-hoc on trained models. In contrast, the \Apeiron{} Framework embeds formal verification at the
architectural level: atomicity is not a property to be checked after
training but a constitutive condition of the irreducible observables
(Layer~1).  A multi-prover pipeline (Section~\ref{subsec:formal-verification})
integrates SymPy, Z3, Lean~4 and Coq; it has produced machine-checkable
atomicity certificates for canonical observables, with Z3 serving as the
sound automatic workhorse and Lean/Coq contributing genuine computational
proofs for Boolean atomicity of integers.

\subsection{Non-Anthropocentric and Emergentist AI}
Recent work has begun to question the anthropocentric assumptions underlying mainstream AI. \textit{Embodied AI} and \textit{enactive cognition} \cite{thompson2007} emphasize the role of the agent's own sensorimotor loop in shaping its categories, but they still presuppose a human-designed body and environment. The \textit{Open-ended Learning} paradigm \cite{stanley2019} seeks systems that continually generate novel behaviors, yet the space of possible behaviors is typically constrained by a human-specified reward or novelty metric. The \Apeiron{} Framework radicalizes these efforts by removing human-defined objectives entirely: the system constructs its own categories, temporalities, and causal structures from first principles, without any external reward signal.

\subsection{Precedents: Research Programme Papers in AI}
\label{subsec:precedents}

The present paper is not a report on a completed, empirically validated system. It belongs to a specific genre of scientific communication: the \emph{research programme paper}. Several influential works in artificial intelligence have successfully employed this format to introduce new paradigms through rigorous conceptual and formal foundations, long before full-scale empirical validation was feasible.

\begin{itemize}
    \item \textbf{Schmidhuber's G\"odel Machine} \cite{schmidhuber2007} proposed a self-referential agent capable of rewriting its own code. Although never fully implemented, its formal specification has inspired subsequent research on self-improving AI and remains a touchstone for theoretical discussions of recursive self-modification.
    
    \item \textbf{Hutter's AIXI} \cite{hutter2005} provides a formal theory of optimal reinforcement learning based on Solomonoff induction. Its primary contribution is mathematical: it defines an optimal agent in a fully general sense, even though the agent is incomputable. AIXI has spawned a rich literature on computable approximations and continues to shape theoretical reinforcement learning.
    
    \item \textbf{Friston's Free Energy Principle} \cite{friston2010} began as a theoretical proposal in neuroscience, postulating that biological systems minimize variational free energy. Over more than a decade, this conceptual framework has expanded into a broad research programme spanning neuroscience, psychiatry, and even artificial agent design, demonstrating how a well-articulated theoretical vision can progressively accumulate empirical support.
\end{itemize}

The \Apeiron{} Framework follows in this tradition. Like these works, it aims to establish a coherent conceptual and formal foundation---in this case, for non-anthropocentric knowledge genesis---upon which future theoretical and empirical work can build. Table~\ref{tab:comparison_precedents} summarizes the comparative scope of these seminal research programme papers and highlights the distinctive contribution of \Apeiron{}.

\begin{table}[htbp]
\centering
\caption{Comparison of \Apeiron{} with foundational research programme papers.}
\label{tab:comparison_precedents}
\begin{tabular}{p{1.6cm} p{2.2cm} p{2.2cm} p{2.2cm} p{2.2cm}}
\toprule
\textbf{Dimension} & \textbf{G\"odel Machine} & \textbf{AIXI} & \textbf{Free Energy Principle} & \textbf{Apeiron} \\
\midrule
Core concept & Self-referential code rewriting & Optimal reward-based agent & Variational free energy minimization & Autonomous knowledge genesis from first principles \\
Ontology & Monolithic (single agent) & Fixed (reward signal) & Hierarchical generative models & Layered, with functorial translations (17 layers) \\
Time & External & External & Implicit in dynamics & Endogenous, emergent \\
Ethics & Not addressed & Via reward specification & Via priors & Embedded as architectural primitives (Layers 15--17) \\
Empirical status & Theoretical construct & Incomputable, approximations exist & Neuroscientific & Cluster I implemented and validated; Clusters II--IV specified \\
\bottomrule
\end{tabular}
\end{table}

\noindent\textbf{Recent developments.} The framework also engages with contemporary advances in open-ended learning \cite{stanley2019}, foundation models, and categorical deep learning \cite{gavranovic2024}. While these approaches push the boundaries of human-designed AI, \Apeiron{} systematically removes the human from the loop of category formation, aiming to construct a truly non-anthropocentric observer.

\subsection{Contrast with Deep Learning Paradigms}
\label{subsec:contrast}

To clarify the distinctive contribution of the \textsc{Apeiron} Framework, Table~\ref{tab:contrast} contrasts its foundational assumptions with those of mainstream deep learning as exemplified by LeCun, Bengio, and Hinton \cite{lecun2015}.

\begin{table}[htbp]
\centering
\caption{Fundamental contrasts between deep learning and the \textsc{Apeiron} Framework.}
\label{tab:contrast}
\begin{tabular}{p{0.22\textwidth} p{0.35\textwidth} p{0.35\textwidth}}
\toprule
\textbf{Aspect} & \textbf{Deep Learning (LeCun et al.)} & \textbf{\textsc{Apeiron} Framework} \\
\midrule
Origin of knowledge & Human-labeled data; supervised or self-supervised objectives & Autonomous emergence from elementary observables; no external reward \\
Role of human & Designs architecture, labels data, defines tasks & Defines only initial atomicity criteria and meta-specification \\
Time & External, absolute step size & Endogenous, emergent from causal relations \\
Explainability & Post-hoc (saliency maps, feature visualization) & Constitutive: traceable to decomposition operators and axioms \\
Ontology & Fixed by human designers & Open-ended; new categories created and reconciled \\
Ethical constraints & External, applied after design & Embedded as architectural primitives (Layers 15--17) \\
\bottomrule
\end{tabular}
\end{table}

This contrast underscores that the \textsc{Apeiron} Framework is not a refinement of existing methods but a fundamentally different paradigm for artificial intelligence.

\subsection{Positioning of the Apeiron Framework}
The \Apeiron{} Framework distinguishes itself through the following unique synthesis:
\begin{itemize}
    \item \textbf{Multi-axial atomicity:} atomicity is validated concurrently across Boolean, measure-theoretic, categorical, and information-theoretic frameworks (Layer~1).
    \item \textbf{Layered ontology:} the seventeen layers form a self-contained progression from irreducible observables to absolute integration, with formal translation functors between layers.
    \item \textbf{Embedded governance:} ethical reversibility and auditability are architectural primitives (Layers 15--17), not external constraints.
    \item \textbf{Falsifiable grounding:} the framework is accompanied by a reference implementation and a suite of falsifiable predictions (Section~\ref{sec:validation}).
    \item \textbf{Research programme framing:} the work is explicitly presented as an open, collaborative research programme, with a clear delineation of current accomplishments and future directions (Section~\ref{sec:status}).
\end{itemize}